\begin{table}[!htbp]
\centering
\caption{Post-outcome exploratory comparison of early and later postperiods}\label{tab:appendix_population_eras}
\small
\setlength{\tabcolsep}{4pt}
\resizebox{\textwidth}{!}{%
\begin{tabular}{lrrrr}
\toprule
Cell construction & 2023--2024 coefficient [95\% CI] & 2025--2026 coefficient [95\% CI] & Later minus early [95\% CI] & Paired $\mathrm{MDE}_{80}$ \\
\midrule
Official final-weight stocks & $-0.111\;[-0.197,-0.024]$ & $-0.166\;[-0.270,-0.062]$ & $-0.056\;[-0.123,0.012]$ & 0.097 \\
Respondent-equivalent routed counts & $-0.103\;[-0.193,-0.012]$ & $-0.183\;[-0.277,-0.089]$ & $-0.080\;[-0.139,-0.021]$ & 0.085 \\
\bottomrule
\end{tabular}}
\tabnote{This timing comparison was produced after outcomes were opened and is post-outcome exploratory.  Both rows use the corrected 113-month calendar, the same 468 occupations, a joint two-era specification, and 9,999 common occupation multipliers within each paired later-minus-early contrast.  Official final-weight stocks are the paper's employment-stock object.  Respondent-equivalent values are fractional routed-record counts after occupational crosswalking; they are neither distinct survey respondents nor a recovered counterfactual CPS weight vintage.  The official-weight interval does not detect an era difference, whereas the respondent-equivalent interval excludes zero.  That divergence does not identify response composition or show that population-control revisions caused the later estimate.}
\end{table}
